
\documentclass{article}
\usepackage{colortbl}
\usepackage{makecell}
\usepackage{multirow}
\usepackage{supertabular}

\begin{document}

\newcounter{utterance}

\centering \large Interaction Transcript for game `codenames', experiment `risk\_low', episode 0 with Apertus{-}v1.5{-}8B.
\vspace{24pt}

{ \footnotesize  \setcounter{utterance}{1}
\setlength{\tabcolsep}{0pt}
\begin{supertabular}{c@{$\;$}|p{.15\linewidth}@{}p{.15\linewidth}p{.15\linewidth}p{.15\linewidth}p{.15\linewidth}p{.15\linewidth}}
    \# & $\;$A & \multicolumn{4}{c}{Game Master} & $\;\:$B\\
    \hline

    \theutterance \stepcounter{utterance}  
    & & & \multicolumn{2}{p{0.3\linewidth}}{
        \cellcolor[rgb]{0.9,0.9,0.9}{
            \makecell[{{p{\linewidth}}}]{
                \texttt{\tiny{[GM$|$GM]}}
                {'hidden': {'team': ['كائن فضائي', 'بوظة', 'ماموث', 'غابة', 'جسر', 'قوة', 'موسكو', 'سرير', 'روبن'], 'innocent': ['لوحة', 'جني', 'إبهام', 'عقد', 'منفذ', 'شريط', 'أفريقيا', 'حفل موسيقي'], 'opponent': ['ليمون', 'سياج', 'كتكوت', 'الفندق', 'حفرة', 'سيارة', 'ختم', 'حديد'], 'assassin': []}, 'revealed': {'team': {'team': [], 'innocent': [], 'opponent': [], 'assassin': []}, 'opponent': {'team': [], 'innocent': [], 'opponent': [], 'assassin': []}}}
            }
        }
    }
    & & \\ \\

    \theutterance \stepcounter{utterance}  
    & & \multicolumn{4}{p{0.6\linewidth}}{
        \cellcolor[rgb]{0.9,0.9,0.9}{
            \makecell[{{p{\linewidth}}}]{
                \texttt{\tiny{[P1$\langle$GM]}}
                \texttt{قدّم تلميحًا واحدًا مكوّنًا من كلمة واحدة يرتبط ببعض كلمات فريقك من القائمة أدناه. يمكنك اختيار عدد الكلمات التي يرتبط بها تلميحك وأيّها.} \\
\texttt{الهدف هو أن يخمّن شريكك جميع كلمات فريقك بناءً على التلميحات التي تقدّمها له، قبل أن يكشف الفريق الآخر جميع كلماته. إذا خَمَّن شريكك كلمة 'الخصم' أو كلمة 'المشتّت'، تنتهي جولة التخمين الخاصة به ويحصل الفريق المنافس على دور. إذا خَمَّن شريكك كلمة 'القاتل'، تنتهي اللعبة ويخسر فريقك فورًا.} \\
\texttt{كلما ارتبط تلميحك بعدد أكبر من الكلمات، كان بإمكانك الفوز أسرع، لكن قد تصبح الروابط أصعب على شريكك في التخمين، لذا اختر كلمة التلميح بحكمة. يجب أن تكون كلمة التلميح مرتبطة دلاليًا بالكلمات المستهدفة، ولا يمكن أن تكون واحدة من الكلمات في القوائم أو أن تحتوي على أجزاء منها.} \\
\\ 
\texttt{قدّم دائمًا تلميحك المكوّن من كلمة واحدة وقائمتك المفصولة بفواصل من الكلمات المستهدفة المرتبطة بالتنسيق التالي واجعل إجاباتك قصيرة قدر الإمكان، ولا تُدرج أي نص آخر غير المطلوب في هذا النموذج:} \\
\\ 
\texttt{تلميح: <كلمة>} \\
\texttt{أهداف: <أهداف>} \\
\\ 
\texttt{كلمات فريقك هي: كائن فضائي، بوظة، ماموث، غابة، جسر، قوة، موسكو، سرير، روبن.} \\
\texttt{كلمات فريق خصمك هي: ليمون، سياج، كتكوت، الفندق، حفرة، سيارة، ختم، حديد.} \\
\texttt{كلمات التشتيت هي: لوحة، جني، إبهام، عقد، منفذ، شريط، أفريقيا، حفل موسيقي.} \\
\texttt{كلمات القتل هي: .} \\
            }
        }
    }
    & \\ \\

    \theutterance \stepcounter{utterance}  
    & \multicolumn{4}{p{0.6\linewidth}}{
        \cellcolor[rgb]{0.9,0.9,0.9}{
            \makecell[{{p{\linewidth}}}]{
                \texttt{\tiny{[P1$\rangle$GM]}}
                \texttt{تلميح: زاوية} \\
\texttt{أهداف: كائن فضائي، بوظة، ماموث، غابة، جسر، قوة، موسكو، سرير، روبن} \\
            }
        }
    }
    & & \\ \\

    \theutterance \stepcounter{utterance}  
    & & & \multicolumn{2}{p{0.3\linewidth}}{
        \cellcolor[rgb]{0.9,0.9,0.9}{
            \makecell[{{p{\linewidth}}}]{
                \texttt{\tiny{[GM$|$GM]}}
                \texttt{زاوية} \\
            }
        }
    }
    & & \\ \\

    \theutterance \stepcounter{utterance}  
    & & & \multicolumn{2}{p{0.3\linewidth}}{
        \cellcolor[rgb]{0.9,0.9,0.9}{
            \makecell[{{p{\linewidth}}}]{
                \texttt{\tiny{[GM$|$GM]}}
                ['كائن فضائي', 'بوظة', 'ماموث', 'غابة', 'جسر', 'قوة', 'موسكو', 'سرير', 'روبن']
            }
        }
    }
    & & \\ \\

    \theutterance \stepcounter{utterance}  
    & & & \multicolumn{2}{p{0.3\linewidth}}{
        \cellcolor[rgb]{0.9,0.9,0.9}{
            \makecell[{{p{\linewidth}}}]{
                \texttt{\tiny{[GM$|$GM]}}
                {'word': 'كائن فضائي', 'assignment': 'team'}
            }
        }
    }
    & & \\ \\

    \theutterance \stepcounter{utterance}  
    & & & \multicolumn{2}{p{0.3\linewidth}}{
        \cellcolor[rgb]{0.9,0.9,0.9}{
            \makecell[{{p{\linewidth}}}]{
                \texttt{\tiny{[GM$|$GM]}}
                {'word': 'بوظة', 'assignment': 'team'}
            }
        }
    }
    & & \\ \\

    \theutterance \stepcounter{utterance}  
    & & & \multicolumn{2}{p{0.3\linewidth}}{
        \cellcolor[rgb]{0.9,0.9,0.9}{
            \makecell[{{p{\linewidth}}}]{
                \texttt{\tiny{[GM$|$GM]}}
                {'word': 'ماموث', 'assignment': 'team'}
            }
        }
    }
    & & \\ \\

    \theutterance \stepcounter{utterance}  
    & & & \multicolumn{2}{p{0.3\linewidth}}{
        \cellcolor[rgb]{0.9,0.9,0.9}{
            \makecell[{{p{\linewidth}}}]{
                \texttt{\tiny{[GM$|$GM]}}
                {'word': 'غابة', 'assignment': 'team'}
            }
        }
    }
    & & \\ \\

    \theutterance \stepcounter{utterance}  
    & & & \multicolumn{2}{p{0.3\linewidth}}{
        \cellcolor[rgb]{0.9,0.9,0.9}{
            \makecell[{{p{\linewidth}}}]{
                \texttt{\tiny{[GM$|$GM]}}
                {'word': 'جسر', 'assignment': 'team'}
            }
        }
    }
    & & \\ \\

    \theutterance \stepcounter{utterance}  
    & & & \multicolumn{2}{p{0.3\linewidth}}{
        \cellcolor[rgb]{0.9,0.9,0.9}{
            \makecell[{{p{\linewidth}}}]{
                \texttt{\tiny{[GM$|$GM]}}
                {'word': 'قوة', 'assignment': 'team'}
            }
        }
    }
    & & \\ \\

    \theutterance \stepcounter{utterance}  
    & & & \multicolumn{2}{p{0.3\linewidth}}{
        \cellcolor[rgb]{0.9,0.9,0.9}{
            \makecell[{{p{\linewidth}}}]{
                \texttt{\tiny{[GM$|$GM]}}
                {'word': 'موسكو', 'assignment': 'team'}
            }
        }
    }
    & & \\ \\

    \theutterance \stepcounter{utterance}  
    & & & \multicolumn{2}{p{0.3\linewidth}}{
        \cellcolor[rgb]{0.9,0.9,0.9}{
            \makecell[{{p{\linewidth}}}]{
                \texttt{\tiny{[GM$|$GM]}}
                {'word': 'سرير', 'assignment': 'team'}
            }
        }
    }
    & & \\ \\

    \theutterance \stepcounter{utterance}  
    & & & \multicolumn{2}{p{0.3\linewidth}}{
        \cellcolor[rgb]{0.9,0.9,0.9}{
            \makecell[{{p{\linewidth}}}]{
                \texttt{\tiny{[GM$|$GM]}}
                {'word': 'روبن', 'assignment': 'team'}
            }
        }
    }
    & & \\ \\

    \theutterance \stepcounter{utterance}  
    & & \multicolumn{4}{p{0.6\linewidth}}{
        \cellcolor[rgb]{0.9,0.9,0.9}{
            \makecell[{{p{\linewidth}}}]{
                \texttt{\tiny{[GM$\rangle$P2]}}
                \texttt{قدّم قائمة مفصولة بفواصل تضم ما يصل إلى 9 من الكلمات من القائمة التالية التي ترتبط على أفضل وجه أو تكون الأكثر ارتباطًا بالكلمة 'زاوية'. ابدأ دائمًا قائمة تخميناتك بـ تخمين: ولا تُدرج أي نص آخر في إجابتك.} \\
\\ 
\texttt{سرير، غابة، جسر، ماموث، قوة، سيارة، شريط، جني، منفذ، روبن، الفندق، حديد، حفل موسيقي، كتكوت، بوظة، إبهام، سياج، عقد، أفريقيا، موسكو، ختم، ليمون، لوحة، كائن فضائي، حفرة} \\
            }
        }
    }
    & \\ \\

    \theutterance \stepcounter{utterance}  
    & & & \multicolumn{4}{p{0.6\linewidth}}{
        \cellcolor[rgb]{0.9,0.9,0.9}{
            \makecell[{{p{\linewidth}}}]{
                \texttt{\tiny{[GM$\langle$P2]}}
                \texttt{تخمين: زاوية, جدار, سقف, أرض, مبنى, مبنى, جدران, زاوية, هيكل} \\
            }
        }
    }
     \\ \\

    \theutterance \stepcounter{utterance}  
    & & & \multicolumn{2}{p{0.3\linewidth}}{
        \cellcolor[rgb]{0.9,0.9,0.9}{
            \makecell[{{p{\linewidth}}}]{
                \texttt{\tiny{[GM$|$GM]}}
                {'player': 'Guesser', 'type': 'guess is clue word', 'utterance': 'زاوية, جدار, سقف, أرض, مبنى, مبنى, جدران, زاوية, هيكل', 'message': "Guessed word 'زاوية' is the same word as the provided clue word, you should only select words from the provided list.", 'guess': 'زاوية', 'clue': 'زاوية'}
            }
        }
    }
    & & \\ \\

\end{supertabular}
}

\end{document}
